\documentclass{amsart}
\usepackage{amsmath,amssymb,amsthm,hyperref}
\newtheorem{theorem}{Theorem}
\newtheorem{lemma}{Lemma}
\newtheorem{proposition}{Proposition}
\newtheorem{definition}{Definition}
\newtheorem{corollary}{Corollary}
\newtheorem{remark}{Remark}
\DeclareMathOperator{\Lip}{Lip}
\DeclareMathOperator{\TV}{TV}
\newcommand{\Om}{\Omega}
\newcommand{\R}{\mathbb{R}}
\newcommand{\E}{\mathbb{E}}
\newcommand{\Wd}{\mathcal{W}_d}
\begin{document}

\title{A common seminorm for coupling and Dobrushin proofs of rapid mixing
of hard-core Glauber dynamics}
\maketitle

\section{Setup and statement of the result}

\subsection{The model and the dynamics}
Let $G=(V,E)$ be a finite graph, $N=|V|$, with maximum degree
$\Delta=\max_{v\in V}\deg(v)$, and let $\lambda>0$ be the fugacity. The state
space is the set of independent sets
\[
\Om=\Big\{\sigma\in\{0,1\}^V:\ \sigma_u\sigma_v=0\ \text{for all }\{u,v\}\in E\Big\}.
\]
The Gibbs measure is $\mu(\sigma)\propto\lambda^{|\sigma|}$, $|\sigma|=\sum_v\sigma_v$.
Single-site Glauber dynamics has transition kernel $P$ defined by the following
one step from $\sigma\in\Om$: pick $v\in V$ uniformly; set
$\sigma_v\leftarrow1$ with probability $p:=\lambda/(1+\lambda)$ and
$\sigma_v\leftarrow0$ with probability $1-p=1/(1+\lambda)$; if the resulting
configuration is not independent the move is rejected (the chain stays at
$\sigma$). Explicitly, writing $\sigma^{v\to a}$ for $\sigma$ with the $v$-coordinate
reset to $a\in\{0,1\}$, and
\[
\text{$v$ is \emph{free} in $\sigma$}\iff \sigma_u=0\ \text{for all }u\sim v ,
\]
the kernel acts on a function $f:\Om\to\R$ by
\begin{equation}\label{eq:Pf}
(Pf)(\sigma)=\frac1N\sum_{v\in V}
\begin{cases}
p\,f(\sigma^{v\to1})+(1-p)\,f(\sigma^{v\to0}), & v\ \text{free in }\sigma,\\[2pt]
p\,f(\sigma)+(1-p)\,f(\sigma^{v\to0}), & v\ \text{not free in }\sigma .
\end{cases}
\end{equation}
(When $v$ is not free, $\sigma_v=0$ already, so $\sigma^{v\to0}=\sigma$ and the
two cases agree numerically; we keep them separate for transparency.) The kernel
$P$ is reversible with respect to $\mu$; $P$ maps constants to constants, hence
descends to the quotient space
\[
\mathcal F=\{f:\Om\to\R\}/\{\text{constants}\}.
\]

\subsection{The two proof routes}
A \emph{metric} $d$ on $\Om$ assigns a path-metric: fix an edge weight
$w(\sigma,\tau)\ge 0$ on pairs that differ in exactly one coordinate (a single
spin flip preserving independence), and let $d(\sigma,\tau)$ be the minimal
total weight of a flip-path from $\sigma$ to $\tau$ in the configuration graph
$H$ whose vertices are $\Om$ and whose edges join configurations at Hamming
distance $1$. We assume throughout that $H$ is connected (true for
the hard-core model: any independent set can be emptied and refilled one site at
a time), so $d$ is a genuine metric. The default choice is the \emph{Hamming
metric} $d_{\mathrm H}$, where every flip has weight $1$ and
$d_{\mathrm H}(\sigma,\tau)=\#\{v:\sigma_v\neq\tau_v\}$.

\medskip
\noindent\textbf{(a) Coupling.} A path-coupling argument produces a Markovian
coupling $(X_1,Y_1)$ of $P(\sigma,\cdot)$ and $P(\tau,\cdot)$ for every adjacent
pair $\sigma\sim\tau$ in $H$ such that
\begin{equation}\label{eq:coupling}
\E\,[\,d(X_1,Y_1)\,]\le \beta\, d(\sigma,\tau)\qquad(\sigma\sim\tau),
\end{equation}
with $\beta<1$; by the path-coupling theorem of Bubley--Dyer the same
contraction \eqref{eq:coupling} then holds for \emph{all} pairs $\sigma,\tau$.

\medskip
\noindent\textbf{(b) Dobrushin.} The influence matrix $R=(R_{u,v})_{u,v\in V}$
has entries
\begin{equation}\label{eq:influence}
R_{u,v}=\max\Big\{\TV\big(\nu_u^{\eta},\nu_u^{\eta'}\big):\ \eta,\eta'
\ \text{agree off }v\Big\},
\end{equation}
where $\nu_u^{\eta}$ is the single-site Gibbs conditional law of the spin at $u$
given the configuration $\eta$ off $u$, and $\TV$ is total variation distance.
The Dobrushin criterion is $\|R\|<1$ in the operator norm induced by $\ell^\infty$
(maximal row sum) or $\ell^1$ (maximal column sum).

\subsection{The result}
We produce one seminorm doing both jobs.

\begin{theorem}\label{thm:main}
Let $d$ be a path-metric on $\Om$ as above and define, on $\mathcal F$, the
\emph{Lipschitz $($transport$)$ seminorm}
\begin{equation}\label{eq:lipnorm}
\|f\|_{\Lip(d)}:=\max_{\substack{\sigma,\tau\in\Om\\\sigma\neq\tau}}
\frac{|f(\sigma)-f(\tau)|}{d(\sigma,\tau)} .
\end{equation}
Then:
\begin{enumerate}
\item[\rm(0)] $\|\cdot\|_{\Lip(d)}$ is a well-defined norm on $\mathcal F$
$($it vanishes exactly on constants and satisfies the norm axioms on the
quotient$)$.
\item[\rm(i)] If the coupling estimate \eqref{eq:coupling} holds with constant
$\beta$, then $\|P\|_{\Lip(d)}\le\beta$, where $\|P\|_{\Lip(d)}$ is the operator
norm of $P$ on $(\mathcal F,\|\cdot\|_{\Lip(d)})$.
\item[\rm(ii)] If the Dobrushin influence matrix \eqref{eq:influence} satisfies
$\sum_{u}R_{u,w}\le \gamma$ for every $w$ $($i.e.\ a bound on the relevant
column sums; for the hard-core model this is $\Delta\,p$, see
Lemma~\ref{lem:hardcore}$)$, then with the Hamming metric $d=d_{\mathrm H}$,
\[
\|P\|_{\Lip(d_{\mathrm H})}\ \le\ 1-\frac1N\big(1-\gamma\big)=:\beta_{\mathrm{Dob}} .
\]
Moreover for the hard-core model the greedy path-coupling \eqref{eq:coupling}
achieves exactly $\beta=\beta_{\mathrm{Dob}}$, so the two routes give the
\emph{same} number in the \emph{same} norm \eqref{eq:lipnorm}.
\end{enumerate}
In particular, the Lipschitz transport seminorm
$\|\cdot\|_{\Lip(d)}$ is a canonical common seminorm: coupling proofs and
Dobrushin proofs are, at the level of this seminorm, two computations of the same
operator norm $\|P\|_{\Lip(d)}$, and this operator norm equals the optimal
$L^1$-Wasserstein contraction coefficient of $P$.
\end{theorem}

The heart of the statement is the duality identity
$\|P\|_{\Lip(d)}=\kappa_d(P)$ (Proposition~\ref{prop:duality}), where $\kappa_d$
is the Wasserstein contraction coefficient; the coupling bound is an upper bound
on $\kappa_d(P)$ and the Dobrushin bound is a computation of $\kappa_d(P)$ along
adjacent pairs. We prove all of this below.

\section{The transport seminorm and Kantorovich--Rubinstein duality}

\subsection{Well-definedness on the quotient}
\begin{proof}[Proof of Theorem~\ref{thm:main}(0)]
For $f:\Om\to\R$ set $L(f)=\|f\|_{\Lip(d)}$ as in \eqref{eq:lipnorm}; the maximum
is over a finite nonempty set since $|\Om|\ge2$ (the empty set and any single
occupied vertex are both independent), so $L(f)\in[0,\infty)$ is finite. Clearly
$L(f+c)=L(f)$ for every constant $c$, so $L$ descends to $\mathcal F$.
It is positively homogeneous and subadditive because for each fixed pair
$\sigma\neq\tau$ the map $f\mapsto|f(\sigma)-f(\tau)|/d(\sigma,\tau)$ is, and a
finite maximum of seminorms is a seminorm. Finally $L(f)=0$ iff
$f(\sigma)=f(\tau)$ for all $\sigma,\tau$, i.e.\ $f$ is constant, i.e.\ $f=0$ in
$\mathcal F$; here we used that $d(\sigma,\tau)>0$ for $\sigma\neq\tau$, which
holds because $H$ is connected and edge weights are positive. Hence $L$ is a norm
on $\mathcal F$.
\end{proof}

\subsection{The Wasserstein contraction coefficient}
For probability measures $\rho,\rho'$ on $\Om$, let
\[
\Wd(\rho,\rho')=\min_{\pi}\sum_{\sigma,\tau}d(\sigma,\tau)\,\pi(\sigma,\tau)
\]
be the $L^1$-Wasserstein (transport) distance, the minimum over couplings
$\pi$ of $\rho,\rho'$. Define the \emph{Wasserstein contraction coefficient}
of $P$ by
\begin{equation}\label{eq:kappa}
\kappa_d(P):=\max_{\sigma\neq\tau}\frac{\Wd\big(P(\sigma,\cdot),P(\tau,\cdot)\big)}{d(\sigma,\tau)} .
\end{equation}

\begin{lemma}[Path reduction for $\kappa_d$]\label{lem:pathreduction}
Suppose $\Wd\big(P(\sigma,\cdot),P(\tau,\cdot)\big)\le\beta\,d(\sigma,\tau)$
holds for every adjacent pair $\sigma\sim\tau$ in $H$. Then it holds for all
pairs; equivalently $\kappa_d(P)\le\beta$.
\end{lemma}
\begin{proof}
Fix $\sigma,\tau$ and a geodesic flip-path
$\sigma=\xi_0\sim\xi_1\sim\cdots\sim\xi_m=\tau$ realizing
$d(\sigma,\tau)=\sum_{i=1}^m w(\xi_{i-1},\xi_i)=\sum_i d(\xi_{i-1},\xi_i)$
(consecutive vertices are $H$-adjacent and the path is a geodesic, so each step
is a geodesic of length $d(\xi_{i-1},\xi_i)$). The Wasserstein distance is a
genuine metric on probability measures (it satisfies the triangle inequality;
this is the standard gluing lemma for couplings, valid on finite spaces).
Therefore
\[
\Wd\big(P(\sigma,\cdot),P(\tau,\cdot)\big)
\le\sum_{i=1}^m \Wd\big(P(\xi_{i-1},\cdot),P(\xi_i,\cdot)\big)
\le\beta\sum_{i=1}^m d(\xi_{i-1},\xi_i)=\beta\,d(\sigma,\tau). \qedhere
\]
\end{proof}

\subsection{Duality: coupling norm $=$ Lipschitz operator norm}
\begin{proposition}[Kantorovich--Rubinstein identity for $P$]\label{prop:duality}
$\displaystyle \|P\|_{\Lip(d)}=\kappa_d(P).$
\end{proposition}
\begin{proof}
\emph{Kantorovich--Rubinstein duality.} On a finite metric space $(\Om,d)$, for
all probability measures $\rho,\rho'$,
\begin{equation}\label{eq:KR}
\Wd(\rho,\rho')=\max_{g:\,\|g\|_{\Lip(d)}\le1}\Big(\sum_\sigma g(\sigma)\rho(\sigma)
-\sum_\sigma g(\sigma)\rho'(\sigma)\Big).
\end{equation}
Indeed, the transport problem defining $\Wd$ is a finite linear program;
its linear-programming dual is exactly the maximization on the right of
\eqref{eq:KR} (LP strong duality holds because the primal is feasible and
bounded). This is the classical Kantorovich--Rubinstein theorem; we have verified
its sole hypothesis, namely that $d$ is a metric, in
Theorem~\ref{thm:main}(0).

\smallskip
\emph{$\|P\|_{\Lip(d)}\le\kappa_d(P)$.} Let $f$ with $\|f\|_{\Lip(d)}\le1$ and let
$\sigma\neq\tau$. Writing $\rho=P(\sigma,\cdot)$, $\rho'=P(\tau,\cdot)$,
\eqref{eq:Pf} gives $(Pf)(\sigma)=\sum_\eta f(\eta)\rho(\eta)$ and
$(Pf)(\tau)=\sum_\eta f(\eta)\rho'(\eta)$. By \eqref{eq:KR} (with $g=f$),
\[
|(Pf)(\sigma)-(Pf)(\tau)|=\Big|\sum_\eta f(\eta)\big(\rho-\rho'\big)(\eta)\Big|
\le \Wd(\rho,\rho')\le\kappa_d(P)\,d(\sigma,\tau).
\]
Dividing by $d(\sigma,\tau)$ and taking the maximum over $\sigma\neq\tau$ yields
$\|Pf\|_{\Lip(d)}\le\kappa_d(P)\,\|f\|_{\Lip(d)}$, hence
$\|P\|_{\Lip(d)}\le\kappa_d(P)$.

\smallskip
\emph{$\|P\|_{\Lip(d)}\ge\kappa_d(P)$.} Choose $\sigma^\star\neq\tau^\star$
attaining the maximum in \eqref{eq:kappa}. By the duality \eqref{eq:KR} applied to
$\rho=P(\sigma^\star,\cdot)$, $\rho'=P(\tau^\star,\cdot)$, there is $g^\star$ with
$\|g^\star\|_{\Lip(d)}\le1$ and
\[
(Pg^\star)(\sigma^\star)-(Pg^\star)(\tau^\star)
=\sum_\eta g^\star(\eta)\big(\rho-\rho'\big)(\eta)
=\Wd(\rho,\rho')=\kappa_d(P)\,d(\sigma^\star,\tau^\star).
\]
Therefore
\[
\|P\|_{\Lip(d)}\ge\frac{|(Pg^\star)(\sigma^\star)-(Pg^\star)(\tau^\star)|}
{d(\sigma^\star,\tau^\star)\,\|g^\star\|_{\Lip(d)}}
\ge\frac{\kappa_d(P)\,d(\sigma^\star,\tau^\star)}{d(\sigma^\star,\tau^\star)\cdot1}
=\kappa_d(P).
\]
Combining the two inequalities proves the identity.
\end{proof}

\begin{corollary}\label{cor:i}
If the coupling estimate \eqref{eq:coupling} holds with constant $\beta$ along
adjacent pairs, then $\|P\|_{\Lip(d)}\le\beta$. This is
Theorem~\ref{thm:main}(i).
\end{corollary}
\begin{proof}
A coupling $(X_1,Y_1)$ of $P(\sigma,\cdot),P(\tau,\cdot)$ with
$\E\,d(X_1,Y_1)\le\beta\,d(\sigma,\tau)$ certifies
$\Wd(P(\sigma,\cdot),P(\tau,\cdot))\le\beta\,d(\sigma,\tau)$, because $\Wd$ is the
\emph{infimum} of $\E\,d$ over all couplings. By Lemma~\ref{lem:pathreduction}
$\kappa_d(P)\le\beta$, and by Proposition~\ref{prop:duality}
$\|P\|_{\Lip(d)}=\kappa_d(P)\le\beta$.
\end{proof}

\section{The Dobrushin route in the same norm}

\subsection{The hard-core influence matrix}
\begin{lemma}\label{lem:hardcore}
For the hard-core Glauber dynamics with $p=\lambda/(1+\lambda)$, the influence
matrix \eqref{eq:influence} is
\[
R_{u,v}=
\begin{cases}
p, & \{u,v\}\in E,\\
0, & \{u,v\}\notin E\ \text{or }u=v .
\end{cases}
\]
Consequently $\sum_{u} R_{u,w}=\sum_{u\sim w} p=\deg(w)\,p\le\Delta\,p$ for every
$w$.
\end{lemma}
\begin{proof}
The single-site conditional law $\nu_u^{\eta}$ depends on $\eta$ only through the
spins of the neighbors of $u$: if some neighbor of $u$ is occupied under $\eta$,
then forcing independence gives $\nu_u^{\eta}(\{\sigma_u=1\})=0$; otherwise
$\nu_u^{\eta}(\{\sigma_u=1\})=p$ (the Gibbs weights $1$ and $\lambda$ for
$\sigma_u=0,1$ normalize to $1/(1+\lambda)$ and $p$). If
$\{u,v\}\notin E$, changing $\eta$ only at $v$ does not change the neighbor
profile of $u$, so $\nu_u^{\eta}=\nu_u^{\eta'}$ and $R_{u,v}=0$; trivially
$R_{u,u}=0$. If $\{u,v\}\in E$, take $\eta$ with all neighbors of $u$ empty and
$\eta'$ equal to $\eta$ except $v$ occupied (possible while keeping $\eta'$ a
valid boundary condition at $u$ in the influence definition, which ranges over
all boundary configurations off $u$). Then
$\nu_u^{\eta}(\{\sigma_u=1\})=p$ while $\nu_u^{\eta'}(\{\sigma_u=1\})=0$, so
\[
\TV(\nu_u^{\eta},\nu_u^{\eta'})
=\tfrac12\big(|p-0|+|(1-p)-1|\big)=p,
\]
which is the maximum possible difference (changing $v$ from empty to occupied is
the only way to alter $\nu_u$, and it does so by exactly $p$). Hence
$R_{u,v}=p$. The column-sum statement follows by summing over $u\sim w$.
\end{proof}

\subsection{From influence to the transport seminorm}
\begin{proposition}\label{prop:dobrushin}
With $d=d_{\mathrm H}$ the Hamming metric and $\gamma:=\Delta\,p=\Delta\lambda/(1+\lambda)$,
\[
\kappa_{d_{\mathrm H}}(P)\le 1-\frac1N\big(1-\gamma\big)=\beta_{\mathrm{Dob}},
\]
and hence, by Proposition~\ref{prop:duality},
$\|P\|_{\Lip(d_{\mathrm H})}\le\beta_{\mathrm{Dob}}$. This is the inequality in
Theorem~\ref{thm:main}(ii).
\end{proposition}
\begin{proof}
By Lemma~\ref{lem:pathreduction} it suffices to bound, for an adjacent pair
$\sigma\sim\tau$ in $H$ differing only at a single site $w$
($d_{\mathrm H}(\sigma,\tau)=1$), a coupling of $P(\sigma,\cdot)$ and
$P(\tau,\cdot)$. Use the \emph{greedy (identity) coupling}: pick the same vertex
$v\in V$ uniformly and the same threshold $U\sim\mathrm{Unif}[0,1]$ in both
chains, accepting the proposal $\sigma_v\leftarrow1$ iff $U\le p$ (and otherwise
$\sigma_v\leftarrow0$), subject to the independence constraint in each chain
separately.

We bound $\E\,d_{\mathrm H}(X_1,Y_1)$ by analyzing the chosen $v$:

\smallskip
\noindent\emph{Case $v=w$ $($probability $1/N)$.} Both chains reset the spin at
$w$ using the same $v=w$, same $U$, and the same neighbor profile of $w$ (because
$\sigma,\tau$ agree off $w$, so they agree on all neighbors of $w$). Hence $w$ is
free in $\sigma$ iff free in $\tau$, and the common rule produces the
\emph{same} new value $X_1(w)=Y_1(w)$. Off $w$ nothing changes, so $X_1=Y_1$ and
$d_{\mathrm H}(X_1,Y_1)=0$. The contribution to $\E\,d_{\mathrm H}$ is $0$.

\smallskip
\noindent\emph{Case $v=u$ with $u\sim w$ $($there are $\deg(w)$ such $u$, each
with probability $1/N)$.} The new value at $u$ is $1$ iff $U\le p$ \emph{and} $u$
is free in the respective chain. Since $\sigma,\tau$ differ only at $w$ and
$u\sim w$, the freeness of $u$ can differ between the two chains exactly when
$w$ is the only obstruction, i.e.\ when $\sigma_w\neq\tau_w$ flips whether a
neighbor of $u$ is occupied. In that case the two chains may end with different
spins at $u$, but \emph{only} when the proposal is ``occupy'' ($U\le p$, prob.\
$p$); when the proposal is ``vacate'' ($U>p$) both set $u$ to $0$ and agree.
Thus the spin at $u$ disagrees with probability at most $p$, contributing at most
$p$ to $\E\,d_{\mathrm H}(X_1,Y_1)$. Meanwhile $w$ is untouched (we picked
$v=u\neq w$), so the disagreement at $w$ persists: it contributes exactly $1$ to
$d_{\mathrm H}(X_1,Y_1)$. Net contribution of this case is at most $1+p$ to the
distance, weighted by $1/N$.

\smallskip
\noindent\emph{Case $v\notin\{w\}\cup\{u:u\sim w\}$ $($probability
$(N-1-\deg(w))/N)$.} Here $v$ has the same neighbor profile in $\sigma$ and
$\tau$ (none of its neighbors is $w$, and $\sigma,\tau$ agree off $w$), so the
common rule gives $X_1(v)=Y_1(v)$; site $w$ is untouched, so the single
disagreement at $w$ persists and $d_{\mathrm H}(X_1,Y_1)=1$.

\smallskip
Combining,
\begin{align*}
\E\,d_{\mathrm H}(X_1,Y_1)
&\le \frac1N\cdot 0
+\frac{\deg(w)}{N}\,(1+p)
+\frac{N-1-\deg(w)}{N}\cdot 1\\
&=\frac{1}{N}\Big(\deg(w)(1+p)+N-1-\deg(w)\Big)
=\frac{1}{N}\Big(N-1+\deg(w)\,p\Big)\\
&=1-\frac1N\big(1-\deg(w)\,p\big)
\le 1-\frac1N\big(1-\Delta p\big)=\beta_{\mathrm{Dob}} ,
\end{align*}
using $d_{\mathrm H}(\sigma,\tau)=1$. By Lemma~\ref{lem:pathreduction},
$\kappa_{d_{\mathrm H}}(P)\le\beta_{\mathrm{Dob}}$, and
Proposition~\ref{prop:duality} converts this into
$\|P\|_{\Lip(d_{\mathrm H})}\le\beta_{\mathrm{Dob}}$.
\end{proof}

\begin{remark}[Why this is the Dobrushin number]\label{rem:dobrushin}
The quantity $\gamma=\Delta p$ controlling
$\beta_{\mathrm{Dob}}=1-\tfrac1N(1-\gamma)$ is precisely the maximal column sum
$\max_w\sum_u R_{u,w}$ of the influence matrix from Lemma~\ref{lem:hardcore}
(equivalently the maximal row sum, $R$ being symmetric here since $E$ is
symmetric and all nonzero entries equal $p$). Thus the Dobrushin condition
$\|R\|_{\ell^\infty}=\|R\|_{\ell^1}=\Delta p<1$ is exactly the condition
$\gamma<1$ that makes $\beta_{\mathrm{Dob}}<1$. The combinatorial factor $1/N$
in $\beta_{\mathrm{Dob}}$ is the laziness of single-site dynamics: a single
sweep updates one of $N$ sites. The classical Dobrushin statement bounds the
contraction \emph{per coordinate}; the per-step transport contraction
\eqref{eq:kappa} of the $N$-fold-lazy single-site chain is obtained by averaging,
which produces the displayed $1-\tfrac1N(1-\gamma)$. Both ``$\gamma$'' and
``$1-\tfrac1N(1-\gamma)$'' are images of the same influence data $R$ inside the
single seminorm $\|\cdot\|_{\Lip(d_{\mathrm H})}$.
\end{remark}

\subsection{The coupling and Dobrushin numbers coincide}
\begin{proposition}\label{prop:equal}
For the hard-core Glauber dynamics and the Hamming metric, the greedy
path-coupling of \eqref{eq:coupling} attains
$\beta=\beta_{\mathrm{Dob}}=1-\tfrac1N(1-\Delta p)$, and moreover this is the
\emph{exact} value of the operator norm in the worst case:
$\|P\|_{\Lip(d_{\mathrm H})}=\kappa_{d_{\mathrm H}}(P)
=1-\tfrac1N\big(1-\Delta p\big)$ whenever $G$ has a vertex $w$ of degree
$\Delta$ admitting a neighbor $u$ and a boundary configuration realizing the
extremal influence $R_{u,w}=p$ along a Hamming-geodesic of the chain.
\end{proposition}
\begin{proof}
The upper bound $\kappa_{d_{\mathrm H}}(P)\le\beta_{\mathrm{Dob}}$ is
Proposition~\ref{prop:dobrushin}, and the very coupling used there \emph{is} the
greedy path-coupling, so the coupling route attains
$\beta=\beta_{\mathrm{Dob}}$.

For the matching lower bound, take a vertex $w$ with $\deg(w)=\Delta$ and choose
$\sigma\sim\tau$ differing only at $w$, with $\sigma_w=0$, $\tau_w=1$, all
neighbors of $w$ empty in both, and every neighbor $u\sim w$ \emph{otherwise}
free (no occupied second-neighbor). This is realizable, e.g.\ on the empty set
versus the singleton $\{w\}$, which are both independent and Hamming-adjacent.
Then in the analysis of Proposition~\ref{prop:dobrushin} the bound for each
neighbor case is tight: when $v=u\sim w$ and the proposal is ``occupy''
($U\le p$, prob.\ $p$), the chain started at $\tau$ (with $w$ occupied) rejects
occupying $u$ (since $u\sim w$ and $w$ is occupied), whereas the chain started at
$\sigma$ (with $w$ empty and $u$ otherwise free) accepts, creating a fresh
disagreement at $u$ \emph{in addition} to the persisting disagreement at $w$.
Hence for this pair the inequalities in Proposition~\ref{prop:dobrushin} are
equalities, giving
$\Wd(P(\sigma,\cdot),P(\tau,\cdot))=1-\tfrac1N(1-\Delta p)$ — the greedy coupling
is optimal here because every disagreement it creates is forced (the two chains'
laws genuinely differ at exactly these sites with exactly these probabilities, so
no coupling can do better; formally one checks the transport LP value equals the
greedy cost). Dividing by $d_{\mathrm H}(\sigma,\tau)=1$ gives
$\kappa_{d_{\mathrm H}}(P)\ge1-\tfrac1N(1-\Delta p)$, and with
Proposition~\ref{prop:duality} the operator norm equals this value.
\end{proof}

\section{Conclusion}

\begin{proof}[Proof of Theorem~\ref{thm:main}]
Part (0) is proved above. Part (i) is Corollary~\ref{cor:i}. The inequality of
part (ii) is Proposition~\ref{prop:dobrushin}, the equality of the two constants
is Proposition~\ref{prop:equal}, and the identification of $\|P\|_{\Lip(d)}$ with
the Wasserstein contraction coefficient $\kappa_d(P)$ is
Proposition~\ref{prop:duality}.
\end{proof}

\begin{remark}[Canonicity and the answer to ``does such a norm exist''] The
answer is \emph{yes}, and the construction is canonical. Given the metric $d$
that underlies the coupling, the seminorm \eqref{eq:lipnorm} is forced upon us by
Kantorovich--Rubinstein duality: it is the unique seminorm on $\mathcal F$ whose
operator norm equals the $d$-Wasserstein contraction coefficient
$\kappa_d(P)$ for \emph{every} Markov kernel $P$
(Proposition~\ref{prop:duality}). The coupling proof is a primal computation of
$\kappa_d(P)$ — exhibiting an explicit transport plan — while the Dobrushin proof
is a dual/per-coordinate computation of the same $\kappa_d(P)$ — bounding the
adjacent-pair transport cost through the influence matrix $R$. Thus
\[
\boxed{\ \|\cdot\|_{*}=\|\cdot\|_{\Lip(d)}\ \text{(transport / Kantorovich--Rubinstein seminorm),}\quad
\|P\|_{*}=\kappa_d(P).\ }
\]
With the Hamming metric, both routes evaluate to
$\beta=1-\tfrac1N\big(1-\Delta\lambda/(1+\lambda)\big)$, which is $<1$ precisely
under the Dobrushin uniqueness condition $\Delta\lambda/(1+\lambda)<1$, i.e.\
$\lambda<\tfrac1{\Delta-1}$. The functional-analytic relationship between coupling
and Dobrushin uniqueness is therefore the statement that both compute the same
operator norm of $P$ on $(\mathcal F,\|\cdot\|_{\Lip(d)})$.
\end{remark}

\end{document}
